\documentclass[12pt]{article}
\usepackage[UTF8]{inputenc}
\usepackage{xeCJK}
\setCJKmainfont{Sarabun}
\usepackage{enumitem}
\usepackage{geometry}
\geometry{a4paper, margin=2.5cm}
\usepackage{setspace}
\onehalfspacing

\begin{document}

\begin{center}
{\Large\textbf{สเปกตรัมและสี}}\\[6pt]
{\normalsize แบบฝึกหัดเพิ่มเติม}\\[6pt]
{\normalsize วิชาฟิสิกส์ ม.ปลาย บทที่ 24}\\[4pt]
{\footnotesize ทั้งหมด 30 ข้อ}
\end{center}
\hrule
\bigskip

\section*{แบบฝึกหัดเพิ่มเติม}
\begin{enumerate}[label=\arabic*., itemsep=0.5\baselineskip, topsep=0.5\baselineskip, leftmargin=*]
\item สีในสเปกตรัมของแสงขาวที่มีความยาวคลื่นยาวที่สุดคือสีใด?
\begin{enumerate}[label=\alph*)., itemsep=0.3\baselineskip, topsep=0.3\baselineskip]
\item ก\ 
\begin{minipage}[t]{0.9\linewidth}
\raggedright แดง
\end{minipage}
\item ข\ 
\begin{minipage}[t]{0.9\linewidth}
\raggedright เหลือง
\end{minipage}
\item ค\ 
\begin{minipage}[t]{0.9\linewidth}
\raggedright เขียว
\end{minipage}
\item ง\ 
\begin{minipage}[t]{0.9\linewidth}
\raggedright น้ำเงิน
\end{minipage}
\end{enumerate}
\vspace{-0.3\baselineskip}
\begin{small}
\textit{คำตอบ: ก}
\end{small}

\item ปรากฏการณ์ที่เห็นรุ้งกินน้ำหลังฝนตกเกิดจากการรวมกันของปรากฏการณ์ใด?
\begin{enumerate}[label=\alph*)., itemsep=0.3\baselineskip, topsep=0.3\baselineskip]
\item ก\ 
\begin{minipage}[t]{0.9\linewidth}
\raggedright การสะท้อนและการหักเห
\end{minipage}
\item ข\ 
\begin{minipage}[t]{0.9\linewidth}
\raggedright การกระจายและการสะท้อนภายใน
\end{minipage}
\item ค\ 
\begin{minipage}[t]{0.9\linewidth}
\raggedright การเลี้ยวเบนและการสะท้อน
\end{minipage}
\item ง\ 
\begin{minipage}[t]{0.9\linewidth}
\raggedright การดูดกลืนและการกระจาย
\end{minipage}
\end{enumerate}
\vspace{-0.3\baselineskip}
\begin{small}
\textit{คำตอบ: ข}
\end{small}

\item รังสีอินฟราเรดใช้ประโยชน์ในด้านใดมากที่สุด?
\begin{enumerate}[label=\alph*)., itemsep=0.3\baselineskip, topsep=0.3\baselineskip]
\item ก\ 
\begin{minipage}[t]{0.9\linewidth}
\raggedright การถ่ายภาพรังสีทางการแพทย์
\end{minipage}
\item ข\ 
\begin{minipage}[t]{0.9\linewidth}
\raggedright การควบคุมรีโมท
\end{minipage}
\item ค\ 
\begin{minipage}[t]{0.9\linewidth}
\raggedright การสื่อสารโทรคมนาคม
\end{minipage}
\item ง\ 
\begin{minipage}[t]{0.9\linewidth}
\raggedright การตรวจวัดอุณหภูมิจากระยะไกล
\end{minipage}
\end{enumerate}
\vspace{-0.3\baselineskip}
\begin{small}
\textit{คำตอบ: ง}
\end{small}

\item เซลล์ประสาทตาแบ่งออกเป็นกี่ชนิด?
\begin{enumerate}[label=\alph*)., itemsep=0.3\baselineskip, topsep=0.3\baselineskip]
\item ก\ 
\begin{minipage}[t]{0.9\linewidth}
\raggedright 1 ชนิด
\end{minipage}
\item ข\ 
\begin{minipage}[t]{0.9\linewidth}
\raggedright 2 ชนิด
\end{minipage}
\item ค\ 
\begin{minipage}[t]{0.9\linewidth}
\raggedright 3 ชนิด
\end{minipage}
\item ง\ 
\begin{minipage}[t]{0.9\linewidth}
\raggedright 4 ชนิด
\end{minipage}
\end{enumerate}
\vspace{-0.3\baselineskip}
\begin{small}
\textit{คำตอบ: ข}
\end{small}

\item การผสมสีแบบ RGB เรียกหลักการอะไร?
\begin{enumerate}[label=\alph*)., itemsep=0.3\baselineskip, topsep=0.3\baselineskip]
\item ก\ 
\begin{minipage}[t]{0.9\linewidth}
\raggedright Subtractive color mixing
\end{minipage}
\item ข\ 
\begin{minipage}[t]{0.9\linewidth}
\raggedright Additive color mixing
\end{minipage}
\item ค\ 
\begin{minipage}[t]{0.9\linewidth}
\raggedright Divisible color mixing
\end{minipage}
\item ง\ 
\begin{minipage}[t]{0.9\linewidth}
\raggedright Primary color mixing
\end{minipage}
\end{enumerate}
\vspace{-0.3\baselineskip}
\begin{small}
\textit{คำตอบ: ข}
\end{small}

\item กฎการแผ่รังสีของวีน์ใช้กับวัตถุประเภทใด?
\begin{enumerate}[label=\alph*)., itemsep=0.3\baselineskip, topsep=0.3\baselineskip]
\item ก\ 
\begin{minipage}[t]{0.9\linewidth}
\raggedright วัตถุที่เป็นแหล่งกำเนิดแสงเท่านั้น
\end{minipage}
\item ข\ 
\begin{minipage}[t]{0.9\linewidth}
\raggedright วัตถุที่เป็น blackbody ทุกชนิด
\end{minipage}
\item ค\ 
\begin{minipage}[t]{0.9\linewidth}
\raggedright วัตถุที่มีอุณหภูมิต่ำกว่า 0°C
\end{minipage}
\item ง\ 
\begin{minipage}[t]{0.9\linewidth}
\raggedright วัตถุที่สะท้อนแสงเท่านั้น
\end{minipage}
\end{enumerate}
\vspace{-0.3\baselineskip}
\begin{small}
\textit{คำตอบ: ข}
\end{small}

\item สเปกตรัมแม่เหล็กไฟฟ้าที่มีความยาวคลื่นยาวกว่าแสงที่ตามองเห็นได้แต่สั้นกว่าไมโครเวฟคือรังสีใด?
\begin{enumerate}[label=\alph*)., itemsep=0.3\baselineskip, topsep=0.3\baselineskip]
\item ก\ 
\begin{minipage}[t]{0.9\linewidth}
\raggedright รังสีเอกซ์
\end{minipage}
\item ข\ 
\begin{minipage}[t]{0.9\linewidth}
\raggedright รังสีอัลตราไวโอเลต
\end{minipage}
\item ค\ 
\begin{minipage}[t]{0.9\linewidth}
\raggedright รังสีอินฟราเรด
\end{minipage}
\item ง\ 
\begin{minipage}[t]{0.9\linewidth}
\raggedright รังสีแกมมา
\end{minipage}
\end{enumerate}
\vspace{-0.3\baselineskip}
\begin{small}
\textit{คำตอบ: ค}
\end{small}

\item ทำไมผลไม้สุกจึงมักมีสีแดงหรือเหลือง?
\begin{enumerate}[label=\alph*)., itemsep=0.3\baselineskip, topsep=0.3\baselineskip]
\item ก\ 
\begin{minipage}[t]{0.9\linewidth}
\raggedright เพราะคลอโรฟิลล์ดูดกลืนแสงเขียว
\end{minipage}
\item ข\ 
\begin{minipage}[t]{0.9\linewidth}
\raggedright เพราะผลไม้สุกดูดกลืนแสงเขียวและน้ำเงิน สะท้อนแสงแดงหรือเหลือง
\end{minipage}
\item ค\ 
\begin{minipage}[t]{0.9\linewidth}
\raggedright เพราะผลไม้สุกเปล่งแสงสีเหล่านั้น
\end{minipage}
\item ง\ 
\begin{minipage}[t]{0.9\linewidth}
\raggedright เพราะแสงแดงทำให้ผลไม้สุก
\end{minipage}
\end{enumerate}
\vspace{-0.3\baselineskip}
\begin{small}
\textit{คำตอบ: ข}
\end{small}

\item รังสีอัลตราไวโอเลต B (UVB) มีช่วงความยาวคลื่นเท่าใด?
\begin{enumerate}[label=\alph*)., itemsep=0.3\baselineskip, topsep=0.3\baselineskip]
\item ก\ 
\begin{minipage}[t]{0.9\linewidth}
\raggedright 100 – 280 nm
\end{minipage}
\item ข\ 
\begin{minipage}[t]{0.9\linewidth}
\raggedright 280 – 315 nm
\end{minipage}
\item ค\ 
\begin{minipage}[t]{0.9\linewidth}
\raggedright 315 – 400 nm
\end{minipage}
\item ง\ 
\begin{minipage}[t]{0.9\linewidth}
\raggedright 400 – 700 nm
\end{minipage}
\end{enumerate}
\vspace{-0.3\baselineskip}
\begin{small}
\textit{คำตอบ: ข}
\end{small}

\item แสงเลเซอร์เป็นแสงที่มีคุณสมบัติพิเศษอย่างไร?
\begin{enumerate}[label=\alph*)., itemsep=0.3\baselineskip, topsep=0.3\baselineskip]
\item ก\ 
\begin{minipage}[t]{0.9\linewidth}
\raggedright มีหลายความยาวคลื่นพร้อมกัน
\end{minipage}
\item ข\ 
\begin{minipage}[t]{0.9\linewidth}
\raggedright มีทิศทางเดียว โครีรีเรนต์สูง ความเข้มสูง
\end{minipage}
\item ค\ 
\begin{minipage}[t]{0.9\linewidth}
\raggedright กระจายตัวอย่างรวดเร็ว
\end{minipage}
\item ง\ 
\begin{minipage}[t]{0.9\linewidth}
\raggedright มีสีขาวเทาเหมือนแสงขาวธรรมดา
\end{minipage}
\end{enumerate}
\vspace{-0.3\baselineskip}
\begin{small}
\textit{คำตอบ: ข}
\end{small}

\item ทำไมเราจึงมองเห็นวัตถุได้?
\begin{enumerate}[label=\alph*)., itemsep=0.3\baselineskip, topsep=0.3\baselineskip]
\item ก\ 
\begin{minipage}[t]{0.9\linewidth}
\raggedright เพราะวัตถุเปล่งแสงออกมา
\end{minipage}
\item ข\ 
\begin{minipage}[t]{0.9\linewidth}
\raggedright เพราะแสงจากแหล่งกำเนิดสะท้อนจากวัตถุเข้าตาเรา
\end{minipage}
\item ค\ 
\begin{minipage}[t]{0.9\linewidth}
\raggedright เพราะแสงถูกดูดกลืนทั้งหมดโดยวัตถุ
\end{minipage}
\item ง\ 
\begin{minipage}[t]{0.9\linewidth}
\raggedright เพราะวัตถุมีสีของตัวเอง
\end{minipage}
\end{enumerate}
\vspace{-0.3\baselineskip}
\begin{small}
\textit{คำตอบ: ข}
\end{small}

\item ในระบบ CMYK ตัวอักษร K ย่อมาจากคำว่าอะไร?
\begin{enumerate}[label=\alph*)., itemsep=0.3\baselineskip, topsep=0.3\baselineskip]
\item ก\ 
\begin{minipage}[t]{0.9\linewidth}
\raggedright Kilo
\end{minipage}
\item ข\ 
\begin{minipage}[t]{0.9\linewidth}
\raggedright Korea
\end{minipage}
\item ค\ 
\begin{minipage}[t]{0.9\linewidth}
\raggedright Key (สีดำเพิ่มความเปรียบต่าง)
\end{minipage}
\item ง\ 
\begin{minipage}[t]{0.9\linewidth}
\raggedright Kernel
\end{minipage}
\end{enumerate}
\vspace{-0.3\baselineskip}
\begin{small}
\textit{คำตอบ: ค}
\end{small}

\item ถ้าวัตถุดำสมบูรณ์มีอุณหภูมิเพิ่มขึ้น 4 เท่า ความยาวคลื่นที่ความเข้มสูงสุดจะเปลี่ยนแปลงอย่างไร?
\begin{enumerate}[label=\alph*)., itemsep=0.3\baselineskip, topsep=0.3\baselineskip]
\item ก\ 
\begin{minipage}[t]{0.9\linewidth}
\raggedright เพิ่มขึ้น 4 เท่า
\end{minipage}
\item ข\ 
\begin{minipage}[t]{0.9\linewidth}
\raggedright ลดลง 4 เท่า
\end{minipage}
\item ค\ 
\begin{minipage}[t]{0.9\linewidth}
\raggedright เพิ่มขึ้น 2 เท่า
\end{minipage}
\item ง\ 
\begin{minipage}[t]{0.9\linewidth}
\raggedright ลดลง 2 เท่า
\end{minipage}
\end{enumerate}
\vspace{-0.3\baselineskip}
\begin{small}
\textit{คำตอบ: ข}
\end{small}

\item รังสีแกมมาเกิดจากกระบวนการใดในนิวเคลียส?
\begin{enumerate}[label=\alph*)., itemsep=0.3\baselineskip, topsep=0.3\baselineskip]
\item ก\ 
\begin{minipage}[t]{0.9\linewidth}
\raggedright การเปลี่ยนแปลงระดับพลังงานอิเล็กตรอน
\end{minipage}
\item ข\ 
\begin{minipage}[t]{0.9\linewidth}
\raggedright การสลายตัวของนิวเคลียส
\end{minipage}
\item ค\ 
\begin{minipage}[t]{0.9\linewidth}
\raggedright การแลกเปลี่ยนอิเล็กตรอนระหว่างอะตอม
\end{minipage}
\item ง\ 
\begin{minipage}[t]{0.9\linewidth}
\raggedright การสั่นของโมเลกุล
\end{minipage}
\end{enumerate}
\vspace{-0.3\baselineskip}
\begin{small}
\textit{คำตอบ: ข}
\end{small}

\item ความยาวคลื่นของแสงสีเขียว (550 nm) เมื่อเทียบกับรังสีไมโครเวฟ (1 cm) จะน้อยกว่ากี่เท่า?
\begin{enumerate}[label=\alph*)., itemsep=0.3\baselineskip, topsep=0.3\baselineskip]
\item ก\ 
\begin{minipage}[t]{0.9\linewidth}
\raggedright 10⁴ เท่า
\end{minipage}
\item ข\ 
\begin{minipage}[t]{0.9\linewidth}
\raggedright 10⁵ เท่า
\end{minipage}
\item ค\ 
\begin{minipage}[t]{0.9\linewidth}
\raggedright 10⁶ เท่า
\end{minipage}
\item ง\ 
\begin{minipage}[t]{0.9\linewidth}
\raggedright 10⁷ เท่า
\end{minipage}
\end{enumerate}
\vspace{-0.3\baselineskip}
\begin{small}
\textit{คำตอบ: ก}
\end{small}

\item ทำไมในเวลากลางคืนตาจึงมองเห็นสีได้ไม่ชัด?
\begin{enumerate}[label=\alph*)., itemsep=0.3\baselineskip, topsep=0.3\baselineskip]
\item ก\ 
\begin{minipage}[t]{0.9\linewidth}
\raggedright เพราะเซลล์กรวยทำงานได้ดีในที่มืด
\end{minipage}
\item ข\ 
\begin{minipage}[t]{0.9\linewidth}
\raggedright เพราะเซลล์แท่งไวต่อแสงน้อยและไม่แยกสีได้
\end{minipage}
\item ค\ 
\begin{minipage}[t]{0.9\linewidth}
\raggedright เพราะเซลล์ประสาทตาตายหมด
\end{minipage}
\item ง\ 
\begin{minipage}[t]{0.9\linewidth}
\raggedright เพราะแสงน้อยทำให้เซลล์กรวยทำงานมากเกินไป
\end{minipage}
\end{enumerate}
\vspace{-0.3\baselineskip}
\begin{small}
\textit{คำตอบ: ข}
\end{small}

\item จอโทรทัศน์แบบเดิม (CRT) ใช้หลักการผสมสีแบบใด?
\begin{enumerate}[label=\alph*)., itemsep=0.3\baselineskip, topsep=0.3\baselineskip]
\item ก\ 
\begin{minipage}[t]{0.9\linewidth}
\raggedright CMYK
\end{minipage}
\item ข\ 
\begin{minipage}[t]{0.9\linewidth}
\raggedright RGB แบบบวก
\end{minipage}
\item ค\ 
\begin{minipage}[t]{0.9\linewidth}
\raggedright Subtractive mixing
\end{minipage}
\item ง\ 
\begin{minipage}[t]{0.9\linewidth}
\raggedright Wavelength division
\end{minipage}
\end{enumerate}
\vspace{-0.3\baselineskip}
\begin{small}
\textit{คำตอบ: ข}
\end{small}

\item ดาวฤกษ์สีแดงมีอุณหภูมิพื้นผิวประมาณเท่าใด?
\begin{enumerate}[label=\alph*)., itemsep=0.3\baselineskip, topsep=0.3\baselineskip]
\item ก\ 
\begin{minipage}[t]{0.9\linewidth}
\raggedright 3000 – 4000 K
\end{minipage}
\item ข\ 
\begin{minipage}[t]{0.9\linewidth}
\raggedright 5000 – 6000 K
\end{minipage}
\item ค\ 
\begin{minipage}[t]{0.9\linewidth}
\raggedright 8000 – 10000 K
\end{minipage}
\item ง\ 
\begin{minipage}[t]{0.9\linewidth}
\raggedright 15000 – 20000 K
\end{minipage}
\end{enumerate}
\vspace{-0.3\baselineskip}
\begin{small}
\textit{คำตอบ: ก}
\end{small}

\item การถ่ายภาพรังสีทางการแพทย์ (X-ray imaging) ใช้รังสีเอกซ์เพราะเหตุใด?
\begin{enumerate}[label=\alph*)., itemsep=0.3\baselineskip, topsep=0.3\baselineskip]
\item ก\ 
\begin{minipage}[t]{0.9\linewidth}
\raggedright รังสีเอกซ์ถูกดูดกลืนโดยกระดูกได้ดีกว่าเนื้อเยื่อ
\end{minipage}
\item ข\ 
\begin{minipage}[t]{0.9\linewidth}
\raggedright รังสีเอกซ์ถูกดูดกลืนโดยเนื้อเยื่อได้ดีกว่ากระดูก
\end{minipage}
\item ค\ 
\begin{minipage}[t]{0.9\linewidth}
\raggedright รังสีเอกซ์ทำให้เนื้อเยื่อเติบโต
\end{minipage}
\item ง\ 
\begin{minipage}[t]{0.9\linewidth}
\raggedright รังสีเอกซ์มองเห็นได้ด้วยตาเปล่า
\end{minipage}
\end{enumerate}
\vspace{-0.3\baselineskip}
\begin{small}
\textit{คำตอบ: ก}
\end{small}

\item แสงสีใดในสเปกตรัมที่ตามองเห็นมีพลังงานโฟตอนสูงที่สุด?
\begin{enumerate}[label=\alph*)., itemsep=0.3\baselineskip, topsep=0.3\baselineskip]
\item ก\ 
\begin{minipage}[t]{0.9\linewidth}
\raggedright แดง
\end{minipage}
\item ข\ 
\begin{minipage}[t]{0.9\linewidth}
\raggedright เหลือง
\end{minipage}
\item ค\ 
\begin{minipage}[t]{0.9\linewidth}
\raggedright เขียว
\end{minipage}
\item ง\ 
\begin{minipage}[t]{0.9\linewidth}
\raggedright ม่วง
\end{minipage}
\end{enumerate}
\vspace{-0.3\baselineskip}
\begin{small}
\textit{คำตอบ: ง}
\end{small}

\item วัตถุดำสมบูรณ์ที่อุณหภูมิ 2000 K จะเปล่งรังสีในช่วงใดของสเปกตรัมมากที่สุด?
\begin{enumerate}[label=\alph*)., itemsep=0.3\baselineskip, topsep=0.3\baselineskip]
\item ก\ 
\begin{minipage}[t]{0.9\linewidth}
\raggedright รังสีแกมมา
\end{minipage}
\item ข\ 
\begin{minipage}[t]{0.9\linewidth}
\raggedright รังสีเอกซ์
\end{minipage}
\item ค\ 
\begin{minipage}[t]{0.9\linewidth}
\raggedright อัลตราไวโอเลต
\end{minipage}
\item ง\ 
\begin{minipage}[t]{0.9\linewidth}
\raggedright อินฟราเรด/แสงแดง
\end{minipage}
\end{enumerate}
\vspace{-0.3\baselineskip}
\begin{small}
\textit{คำตอบ: ง}
\end{small}

\item คนที่มองเห็นสีปกติ (trichromats) มีเซลล์กรวยกี่ชนิด และรับแสงสีอะไรได้?
\begin{enumerate}[label=\alph*)., itemsep=0.3\baselineskip, topsep=0.3\baselineskip]
\item ก\ 
\begin{minipage}[t]{0.9\linewidth}
\raggedright 2 ชนิด - แดงและน้ำเงิน
\end{minipage}
\item ข\ 
\begin{minipage}[t]{0.9\linewidth}
\raggedright 3 ชนิด - แดง เขียว น้ำเงิน
\end{minipage}
\item ค\ 
\begin{minipage}[t]{0.9\linewidth}
\raggedright 4 ชนิด - แดง ส้ม เขียว น้ำเงิน
\end{minipage}
\item ง\ 
\begin{minipage}[t]{0.9\linewidth}
\raggedright 3 ชนิด - แดง เหลือง น้ำเงิน
\end{minipage}
\end{enumerate}
\vspace{-0.3\baselineskip}
\begin{small}
\textit{คำตอบ: ข}
\end{small}

\item ถ้าต้องการพิมพ์สีเขียวบนกระดาษขาวโดยใช้ระบบ CMYK ต้องใช้หมึกสีใดผสมกัน?
\begin{enumerate}[label=\alph*)., itemsep=0.3\baselineskip, topsep=0.3\baselineskip]
\item ก\ 
\begin{minipage}[t]{0.9\linewidth}
\raggedright Cyan + Yellow
\end{minipage}
\item ข\ 
\begin{minipage}[t]{0.9\linewidth}
\raggedright Cyan + Magenta
\end{minipage}
\item ค\ 
\begin{minipage}[t]{0.9\linewidth}
\raggedright Magenta + Yellow
\end{minipage}
\item ง\ 
\begin{minipage}[t]{0.9\linewidth}
\raggedright Cyan + Magenta + Yellow
\end{minipage}
\end{enumerate}
\vspace{-0.3\baselineskip}
\begin{small}
\textit{คำตอบ: ก}
\end{small}

\item รังสี UVC มีความยาวคลื่น < 280 nm ใช้ในการฆ่าเชื้อโรค เพราะเหตุใดรังสีนี้จึงอันตรายต่อมนุษย์?
\begin{enumerate}[label=\alph*)., itemsep=0.3\baselineskip, topsep=0.3\baselineskip]
\item ก\ 
\begin{minipage}[t]{0.9\linewidth}
\raggedright เพราะทำให้น้ำเดือด
\end{minipage}
\item ข\ 
\begin{minipage}[t]{0.9\linewidth}
\raggedright เพราะดูดกลืนโดยผิวหนังได้ดีและทำลาย DNA
\end{minipage}
\item ค\ 
\begin{minipage}[t]{0.9\linewidth}
\raggedright เพราะทำให้อุณหภูมิร่างกายสูงขึ้นเร็ว
\end{minipage}
\item ง\ 
\begin{minipage}[t]{0.9\linewidth}
\raggedright เพราะกระเจิงในบรรยากาศไม่ได้
\end{minipage}
\end{enumerate}
\vspace{-0.3\baselineskip}
\begin{small}
\textit{คำตอบ: ข}
\end{small}

\item ถ้าวัตถุ A มี λ\_max = 600 nm และวัตถุ B มี λ\_max = 300 nm อัตราส่วน T\_A / T\_B มีค่าเท่าใด?
\begin{enumerate}[label=\alph*)., itemsep=0.3\baselineskip, topsep=0.3\baselineskip]
\item ก\ 
\begin{minipage}[t]{0.9\linewidth}
\raggedright 1/2
\end{minipage}
\item ข\ 
\begin{minipage}[t]{0.9\linewidth}
\raggedright 1
\end{minipage}
\item ค\ 
\begin{minipage}[t]{0.9\linewidth}
\raggedright 2
\end{minipage}
\item ง\ 
\begin{minipage}[t]{0.9\linewidth}
\raggedright 4
\end{minipage}
\end{enumerate}
\vspace{-0.3\baselineskip}
\begin{small}
\textit{คำตอบ: ก}
\end{small}

\item การเรืองแสงของสาร (fluorescence) ทำงานโดยหลักการใด?
\begin{enumerate}[label=\alph*)., itemsep=0.3\baselineskip, topsep=0.3\baselineskip]
\item ก\ 
\begin{minipage}[t]{0.9\linewidth}
\raggedright สารดูดกลืนแสงแล้วเปล่งแสงที่มีความยาวคลื่นเท่าเดิม
\end{minipage}
\item ข\ 
\begin{minipage}[t]{0.9\linewidth}
\raggedright สารดูดกลืนแสง UV แล้วเปล่งแสงที่ตามองเห็นได้ (ความยาวคลื่นยาวกว่า)
\end{minipage}
\item ค\ 
\begin{minipage}[t]{0.9\linewidth}
\raggedright สารสะท้อนแสงทุกสี
\end{minipage}
\item ง\ 
\begin{minipage}[t]{0.9\linewidth}
\raggedright สารเปลี่ยนแสงให้เป็นคลื่นเสียง
\end{minipage}
\end{enumerate}
\vspace{-0.3\baselineskip}
\begin{small}
\textit{คำตอบ: ข}
\end{small}

\item ทำไมแผ่นกรองแสงโพลารอยด์จึงสามารถลดแสงจ้าได้?
\begin{enumerate}[label=\alph*)., itemsep=0.3\baselineskip, topsep=0.3\baselineskip]
\item ก\ 
\begin{minipage}[t]{0.9\linewidth}
\raggedright เพราะมันดูดกลืนแสงทุกสี
\end{minipage}
\item ข\ 
\begin{minipage}[t]{0.9\linewidth}
\raggedright เพราะมันยอมให้แสงที่สั่นในระนาบเดียวกับแกนของมันผ่านไป
\end{minipage}
\item ค\ 
\begin{minipage}[t]{0.9\linewidth}
\raggedright เพราะมันเปลี่ยนแสงให้เป็นรังสีอินฟราเรด
\end{minipage}
\item ง\ 
\begin{minipage}[t]{0.9\linewidth}
\raggedright เพราะมันกระจายแสงออกไปทุกทิศทาง
\end{minipage}
\end{enumerate}
\vspace{-0.3\baselineskip}
\begin{small}
\textit{คำตอบ: ข}
\end{small}

\item เมื่อใช้แสงสีน้ำเงินส่องลงบนวัตถุสีเหลือง ตาจะเห็นวัตถุนั้นเป็นสีใด?
\begin{enumerate}[label=\alph*)., itemsep=0.3\baselineskip, topsep=0.3\baselineskip]
\item ก\ 
\begin{minipage}[t]{0.9\linewidth}
\raggedright เหลือง
\end{minipage}
\item ข\ 
\begin{minipage}[t]{0.9\linewidth}
\raggedright เขียว
\end{minipage}
\item ค\ 
\begin{minipage}[t]{0.9\linewidth}
\raggedright ดำหรือมืดมาก
\end{minipage}
\item ง\ 
\begin{minipage}[t]{0.9\linewidth}
\raggedright ขาว
\end{minipage}
\end{enumerate}
\vspace{-0.3\baselineskip}
\begin{small}
\textit{คำตอบ: ค}
\end{small}

\item กาแล็กซีทางซ้ายมือในกล้องโทรทรรศน์มีสีแดงกว่าปกติ (redshift) หมายความว่าอย่างไร?
\begin{enumerate}[label=\alph*)., itemsep=0.3\baselineskip, topsep=0.3\baselineskip]
\item ก\ 
\begin{minipage}[t]{0.9\linewidth}
\raggedright กาแล็กซีเผาไหม้กำลังแผ่รังสีสีแดง
\end{minipage}
\item ข\ 
\begin{minipage}[t]{0.9\linewidth}
\raggedright กาแล็กซีกำลังเคลื่อนที่เข้าหาโลก
\end{minipage}
\item ค\ 
\begin{minipage}[t]{0.9\linewidth}
\raggedright กาแล็กซีกำลังเคลื่อนที่ออกจากโลกและสเปกตรัมเลื่อนไปทางแดง
\end{minipage}
\item ง\ 
\begin{minipage}[t]{0.9\linewidth}
\raggedright กาแล็กซีมีอุณหภูมิสูงมาก
\end{minipage}
\end{enumerate}
\vspace{-0.3\baselineskip}
\begin{small}
\textit{คำตอบ: ค}
\end{small}

\item ถ้าเปรียบเทียบพลังงานโฟตอนของรังสีเอกซ์ (λ = 0.1 nm) กับรังสีแกมมา (λ = 0.001 nm) อัตราส่วน E\_X-ray / E\_gamma มีค่าเท่าใด?
\begin{enumerate}[label=\alph*)., itemsep=0.3\baselineskip, topsep=0.3\baselineskip]
\item ก\ 
\begin{minipage}[t]{0.9\linewidth}
\raggedright 0.01
\end{minipage}
\item ข\ 
\begin{minipage}[t]{0.9\linewidth}
\raggedright 0.1
\end{minipage}
\item ค\ 
\begin{minipage}[t]{0.9\linewidth}
\raggedright 10
\end{minipage}
\item ง\ 
\begin{minipage}[t]{0.9\linewidth}
\raggedright 100
\end{minipage}
\end{enumerate}
\vspace{-0.3\baselineskip}
\begin{small}
\textit{คำตอบ: ข}
\end{small}

\end{enumerate}
\end{document}